\documentclass[11pt,a4paper]{article}
\usepackage[margin=2.5cm]{geometry}
\usepackage{amsmath,amssymb}
\usepackage{booktabs}
\usepackage{adjustbox}
\usepackage{enumitem}
\usepackage{xcolor}
\usepackage{xurl}
\usepackage[colorlinks=true,linkcolor=blue,citecolor=blue,urlcolor=blue]{hyperref}

\newcommand{\Esym}{E_{\rm sym}}
\newcommand{\Lsym}{L_{\rm sym}}
\newcommand{\Msun}{M_\odot}
\newcommand{\doi}[1]{\href{https://doi.org/#1}{doi:#1}}
\newcommand{\arxiv}[1]{\href{https://arxiv.org/abs/#1}{arXiv:#1}}

\title{How to proceed: priorities, new results and research directions}
\author{Nisha Singh\\ Shiv Nadar Institution of Eminence}
\date{24 September 2026}

\begin{document}
\sloppy
\maketitle

\section*{In one paragraph}
The three papers are sound and close to submission; what they need is a few days of cleaning, not new physics. The most important change is in how paper~1 is \emph{framed}. My new calculations (Sec.~\ref{sec:what}) show that cooling does not really measure $\Lsym$: it measures the conditions at about twice saturation density, i.e.\ the proton fraction $x_p(2n_0)$, the symmetry energy $S(2n_0)$, and through them the direct Urca threshold mass $M_{\rm DU}$. $\Lsym$ is a model-dependent translation of that. Written this way, the result agrees with Montefusco et al.\ \cite{Montefusco2026} on the physical quantity (the onset mass) and explains the $\Lsym$ tension instead of being weakened by it. The next research step (paper~4) is to run the same cooling analysis with an EoS family in which $\Lsym$ and $S(2n_0)$ are not locked together.

%=====================================================================
\section{Immediate tasks (1--2 weeks)}
%=====================================================================
\begin{enumerate}[leftmargin=*]
\item \textbf{Send me \texttt{Data/cooling\_chi2\_final.csv}} (or run it yourself):
\begin{verbatim}
python3 combine_cooling_lab.py Data/cooling_chi2_final.csv
\end{verbatim}
This gives the credible intervals for $\Esym$, $\Lsym$, $S(2n_0)$, $x_p(2n_0)$, $R_{1.4}$ and $M_{\rm DU}$, the posterior probability $P(M_{\rm DU}\le1.4\,\Msun)$, and the same with the INDRA--FAZIA constraint added (Sec.~\ref{sec:if}). $M_{\rm DU}$ is currently known for 683 of the 864 EoS; run the script on the Mac, where all \texttt{MRcurve\_*\_CAT.dat} files are, after regenerating \texttt{reports/data/eos\_physical\_quantities.dat}, or send me the remaining curves.
\item \textbf{Paper~2:} wait for the neutron-skin rerun (about 10 hours), then regenerate all its numbers (I will do this).
\item \textbf{Fix the citations} marked \texttt{\% CHECK} in the drafts: the RX~J0822 sources (Table~1 of Potekhin et al.\ 2020), the published $\zeta$ values.
\item \textbf{Quick robustness runs} (no new NSCool runs needed): RX~J0822 with $\sigma=0.12$ dex; the $^3P_2$-gap alternative (\texttt{durca\_sf\_altB\_vc}, already run) as a second model in a Bayesian average instead of a separate check.
\item \textbf{Confirm the Landau mass} in the tables actually used (two \texttt{grep} commands in the analysis note).
\item \textbf{Write to D.~Page} about the \texttt{bf\_r} problem before paper~3 is submitted, and ask which NSCool versions are affected.
\end{enumerate}

%=====================================================================
\section{What cooling actually measures}\label{sec:what}
%=====================================================================
For all 864 EoS I computed $S(2n_0)$, $S(3n_0)$, $x_p(2n_0)$, $x_p(3n_0)$ and $K_{\rm sym}$ with the EoS code, and $M_{\rm DU}$ from the pipeline $M$--$R$ curves (\texttt{reports/data/eos\_physical\_quantities.dat}). The correlation of the threshold mass with each quantity:
\begin{center}
\begin{tabular}{lcccccc}
\toprule
 & $\Lsym$ & $\Esym$ & $S(2n_0)$ & $S(3n_0)$ & $x_p(2n_0)$ & $x_p(3n_0)$\\
\midrule
correlation with $M_{\rm DU}$ & $-0.69$ & $-0.72$ & $-0.82$ & $-0.71$ & $\mathbf{-0.85}$ & $-0.76$\\
\bottomrule
\end{tabular}
\end{center}
At the \emph{same} $\Esym=32.5$ and $\Lsym=65$ MeV the five models give:
\begin{center}
\begin{tabular}{lccc}
\toprule
Model & $S(2n_0)$ (MeV) & $x_p(2n_0)$ & $M_{\rm DU}$ ($\Msun$)\\
\midrule
IOPB-I & 53.0 & 0.114 & 1.54\\
FSUGarnet & 53.2 & 0.113 & 1.57\\
BigApple & 54.1 & 0.115 & 2.05\\
GM1 & 49.1 & 0.099 & no DUrca\\
GM2 & 47.2 & 0.093 & no DUrca\\
\bottomrule
\end{tabular}
\end{center}
And conversely, the EoS with $M_{\rm DU}=1.45$--$1.55\,\Msun$ span $\Lsym=50$--85 MeV but only $x_p(2n_0)=0.111$--0.129. So:
\begin{itemize}
\item the cooling data select the proton fraction around $2n_0$ (and hence $M_{\rm DU}$);
\item the $\Lsym$ interval is the image of that selection through five specific Lagrangians, and depends on them;
\item this is why the $\Lsym$ values of cooling and of laboratory analyses can differ while the onset masses agree (1.77 vs 1.89 $\Msun$).
\end{itemize}
\textbf{Suggestion for paper~1:} quote $x_p(2n_0)$, $S(2n_0)$ and $M_{\rm DU}$ as the primary, less model-dependent result, and $\Lsym$ as the model-family translation. This makes the paper stronger, not weaker, against the comparison with \cite{Montefusco2026}.

%=====================================================================
\section{Combining cooling with the heavy-ion constraint}\label{sec:if}
%=====================================================================
Montefusco et al.\ publish, for use by others, a Gaussian form of their INDRA--FAZIA constraint in $(\Esym,\Lsym,K_{\rm sym})$ (their Appendix D). I applied it to our grid:
\begin{itemize}
\item within our model family it prefers $\Esym\approx28$--30 MeV and $\Lsym\approx40$--55 MeV;
\item 28 EoS inside the cooling 68\% region are within $1\sigma$ of it (the closest is GM2 at $\Esym=29.5$, $\Lsym=50$), and 64 within $2\sigma$;
\item our single best fit (GM2 at 32.5/80) is \emph{not}: its $\Esym$ is $\approx3.8\sigma$ above theirs.
\end{itemize}
So cooling and heavy-ion data are compatible, and together they point to $\Esym\approx29$ MeV and $\Lsym\approx50$ MeV. The combination is a natural new result for paper~1 (a paragraph and a figure). It needs only the final cooling $\chi^2$ file; the script \texttt{combine\_cooling\_lab.py} is ready and tested.

Note: the dry run of the script on the old (frozen-tail, one-mode) $\chi^2$ file shows the expected behaviour: adding the heavy-ion constraint pulls $\Esym$ to 28--30 MeV and makes DUrca in $\le1.4\,\Msun$ stars very unlikely. Those numbers are not results; only the final $\chi^2$ file gives real ones.

%=====================================================================
\section{The main limitation, and paper~4}
%=====================================================================
Below $\Lsym=30$ MeV the present model family essentially stops: GM1 gives no usable EoS, GM2 survives at 20--25 MeV but DUrca then never opens, and FSUGarnet, IOPB-I and BigApple end at 35--50 MeV (no real $g_\rho$). Therefore the region preferred by the nuclear-structure part of \cite{Montefusco2026} ($\Lsym\approx20$--37 MeV) cannot be tested with our current models. And in all five models the $\omega$--$\rho$ term ties $\Lsym$ rigidly to $S(2n_0)$.

\textbf{Proposal for paper~4:} repeat the cooling analysis with an EoS family where $\Lsym$ and the high-density symmetry energy are independent. Three options, in order of effort:
\begin{enumerate}[leftmargin=*]
\item \textbf{Density-dependent RMF} (DDME2, DDMEX, DD2 type), where the $\rho$ coupling falls with density as $g_\rho(n)=g_\rho(n_0)\exp[-a_\rho(n/n_0-1)]$. $a_\rho$ controls the high-density behaviour separately from $\Lsym$. This is exactly the parameter Thapa and Sinha tuned for PREX-2 \cite{ThapaSinha2022}, so it links naturally to their work (and is a possible collaboration). The EoS code needs rearrangement terms in the chemical potentials; the rest of the pipeline is unchanged.
\item \textbf{Adding the scalar-isovector $\delta$ meson.} It changes the high-density symmetry energy and splits the proton and neutron Dirac masses, which enters the Landau masses and hence the cooling rates. Recent work already studies its effect on DUrca.
\item \textbf{A metamodel-type EoS} as in \cite{Montefusco2026}, with $\Lsym$ and $K_{\rm sym}$, $Q_{\rm sym}$ free. This would allow a like-for-like comparison with their result, with cooling as the new ingredient.
\end{enumerate}
The key question for paper~4: \emph{with $\Lsym$ and $S(2n_0)$ decoupled, does cooling still prefer $\Lsym\approx50$--75 MeV, or does it only constrain $x_p(2n_0)$?} Either answer is publishable, and the second would reconcile cooling with the laboratory results.

%=====================================================================
\section{Further improvements of the cooling analysis}
%=====================================================================
\begin{enumerate}[leftmargin=*]
\item \textbf{Finer mass grid.} Cooling is run at 1.0, 1.4, 1.8, 2.0 $\Msun$ only, while the TOV profiles already exist every $0.1\,\Msun$. The jump from 1.4 to 1.8 $\Msun$ is exactly where DUrca switches on (median $M_{\rm DU}=1.77\,\Msun$), so the star-by-star mass choice is coarse there. Running 1.1--1.9 $\Msun$ in steps of 0.1 for the final mode costs about $864\times5\times$ a few seconds per run, i.e.\ a few hours on four workers, and allows a continuous mass (with the $1.4\pm0.3\,\Msun$ prior as a proper Bayesian prior).
\item \textbf{Treat each star's mass as a nuisance parameter} and marginalise (sum over masses with the prior) instead of taking the minimum. This is the standard Bayesian choice and removes the discreteness.
\item \textbf{Age uncertainties.} Characteristic spin-down ages can be wrong by factors of a few; kinematic and remnant ages are better. Add an age error per star (larger for characteristic ages) through the local slope of the cooling curve.
\item \textbf{Average over pairing models} (gap ``a'' and ``b'', which are both already run with heating) with equal prior weight, instead of fixing one.
\item \textbf{Use more information on Cas~A.} Besides its temperature, the observed temperature change over two decades has been used to constrain superfluidity (e.g.\ \cite{Shternin2011,Ho2015}). Its current status is debated because of instrumental effects, so it should enter as a systematic check, not a main constraint.
\item \textbf{Report the fit quality honestly per star} (a table of $\chi^2_i$ at the best fit): it shows at once which stars drive the result.
\end{enumerate}

%=====================================================================
\section{Suggested order of publication}
%=====================================================================
\begin{enumerate}[leftmargin=*]
\item \textbf{Paper~1 and paper~3 together} (same data set and pipeline; paper~3 is the robustness check of paper~1 and also reports a code problem that others need to know about). Paper~1 with the reframing of Sec.~\ref{sec:what} and the heavy-ion combination of Sec.~\ref{sec:if}.
\item \textbf{Paper~2} after the skin rerun. Its message (the models give nearly equal Pb and Ca skins; cooling does not choose between PREX-2 and CREX) is robust.
\item \textbf{Paper~4}: density-dependent or $\delta$-meson family through the same pipeline.
\end{enumerate}

\begin{thebibliography}{9}
\small
\bibitem{Montefusco2026} G.~Montefusco et al., \emph{Ruling out nucleonic direct Urca cooling in low-mass neutron stars using nuclear data}, \arxiv{2609.24940} (2026).
\bibitem{ThapaSinha2022} V.~B.~Thapa and M.~Sinha, \emph{Direct URCA process in light of PREX-2}, \arxiv{2203.02272}.
\bibitem{Shternin2011} P.~S.~Shternin, D.~G.~Yakovlev, C.~O.~Heinke, W.~C.~G.~Ho and D.~J.~Patnaude, Mon.\ Not.\ R.\ Astron.\ Soc.\ \textbf{412}, L108 (2011), \doi{10.1111/j.1745-3933.2011.01015.x}.
\bibitem{Ho2015} W.~C.~G.~Ho, K.~G.~Elshamouty, C.~O.~Heinke and A.~Y.~Potekhin, Phys.\ Rev.\ C \textbf{91}, 015806 (2015), \doi{10.1103/PhysRevC.91.015806}.
\end{thebibliography}
\end{document}
